\chapter{Experimental Setup}
\label{chap:experimental-setup}

\section{Evaluation Objectives}

The experiment evaluates the final proposed pipeline rather than training a new
retrieval or generation model. Its purpose is to measure whether the system:

\begin{itemize}
    \item understands the user's current travel request;
    \item distinguishes intent, operation, explicit constraints, and retrieval
    facets;
    \item retrieves and ranks useful evidence;
    \item applies relevant persistent preferences and temporary context;
    \item generates correct, faithful, personalized, and sufficiently complete
    answers; and
    \item operates with observable latency, token usage, and estimated cost.
\end{itemize}

No legacy ablation configuration is treated as the principal system. The
official experiment evaluates the current v2.5 pipeline. Stage-level traces,
recovery-before/after comparisons, and error categories are used to explain
behavior within this final configuration.

\section{Corpus Snapshot}

All evaluated queries use one fixed tourism-corpus snapshot. Before the run, the
experiment records a corpus fingerprint derived from ordered chunk identifiers
and update information, together with:

\begin{itemize}
    \item document count;
    \item chunk count;
    \item number of chunks with embeddings;
    \item embedding-model distribution; and
    \item latest corpus update timestamp.
\end{itemize}

The final values from \texttt{02\_pipeline\_traces.json} must be copied into
Table~\ref{tab:corpus-snapshot} after the official run.

\begin{table}[H]
\centering
\caption{Corpus snapshot used in the official experiment}
\label{tab:corpus-snapshot}
\begin{tabular}{ll}
\hline
\textbf{Field} & \textbf{Official value} \\
\hline
Documents & \textit{Insert from official trace} \\
Chunks & \textit{Insert from official trace} \\
Embedded chunks & \textit{Insert from official trace} \\
Embedding model & \texttt{all-MiniLM-L6-v2} \\
Embedding dimension & 384 \\
Corpus fingerprint & \textit{Insert abbreviated fingerprint} \\
Latest update & \textit{Insert from official trace} \\
\hline
\end{tabular}
\end{table}

\section{Dataset Design}

The evaluation dataset contains five synthetic users. Profiles are deliberately
different in expertise, interests, travel styles, preferred activities, budget,
avoidance preferences, answer length, tone, and explanation style. This
variation provides controlled personalization requirements that are difficult
to obtain from unrelated one-shot queries.

Each user has two conversations, and each conversation contains ten turns. The
complete design is therefore:

\begin{equation}
5\ \text{users} \times 2\ \text{conversations} \times
10\ \text{turns} = 100\ \text{query turns}.
\end{equation}

Queries within a conversation are generated as a coherent sequence. Later
turns may use references, preserve selected destinations, refine constraints,
or ask comparisons based on previous answers. The two conversations for the
same user allow durable profile information to remain relevant while temporary
trip state changes.

Each generated turn contains draft reference annotations for intent, operation,
explicit constraints, retrieval facets, applicable personalization, relevant
memory identifiers, and expected answer facts when applicable. These labels are
not accepted automatically as ground truth. Human review checks taxonomy
consistency, location types, negative constraints, and ambiguous intent labels
before final reporting.

\section{Model and Pipeline Configuration}

Table~\ref{tab:model-roles-experiment} shows the model assignments used in the
current experimental configuration. Model identifiers, pricing information,
and prompt hashes are also stored in the run artifacts.

\begin{table}[H]
\centering
\caption{Model roles in the official experiment}
\label{tab:model-roles-experiment}
\begin{tabular}{p{0.40\textwidth} p{0.45\textwidth}}
\hline
\textbf{Role} & \textbf{Model} \\
\hline
Conversational rewriting & \texttt{deepseek-v4-flash} \\
Structured query parsing & \texttt{deepseek-v4-pro} \\
Planning and evidence checking & \texttt{deepseek-v4-flash} \\
Answer generation and memory & \texttt{deepseek-v4-flash} \\
Synthetic dataset generation & \texttt{gpt-5.6-luna} \\
Independent LLM judge & \texttt{gpt-5.6-terra} \\
Dense embedding & \texttt{all-MiniLM-L6-v2} \\
CrossEncoder reranker & \texttt{cross-encoder/ms-marco-MiniLM-L-6-v2} \\
\hline
\end{tabular}
\end{table}

The principal retrieval parameters are:

\begin{itemize}
    \item up to six planner requirements and six retrieval tasks;
    \item normally 5--20 dense and BM25 candidates per task, with fallback 10;
    \item a pipeline-level ceiling of 30 candidates per retriever;
    \item RRF with $k=60$;
    \item up to 20 fused candidates passed to the reranker;
    \item up to eight selected evidence chunks; and
    \item up to three focused recovery queries.
\end{itemize}

\section{Experiment Stages}

The official experiment is executed as four resumable stages:

\begin{enumerate}
    \item \textbf{Generate}: create profiles, memories, conversations, queries,
    and draft references in \texttt{01\_generated\_dataset.json}.

    \item \textbf{Pipeline}: execute every query through the final RAG pipeline
    and save all intermediate stages in
    \texttt{02\_pipeline\_traces.json}.

    \item \textbf{Judge}: evaluate candidate relevance and final-answer quality
    using the independent judge, saving
    \texttt{03\_judged\_traces.json}.

    \item \textbf{Export}: validate the judged schema, calculate deterministic
    metrics, and produce \texttt{04\_thesis\_dataset.json} and
    \texttt{05\_metric\_summary.json}.
\end{enumerate}

Generation, answering, and judging use separate prompts. The answer model and
judge model are assigned different roles, reducing direct self-evaluation.
This does not eliminate common model biases, so human confirmation remains part
of the protocol.

\section{Query-Understanding Metrics}

\subsection{Intent and Operation Accuracy}

Intent Accuracy is the proportion of cases whose predicted intent matches the
validated reference:

\begin{equation}
\operatorname{IntentAccuracy}
=\frac{1}{N}\sum_{i=1}^{N}
\mathbb{I}(\hat{y}^{\mathrm{intent}}_i=y^{\mathrm{intent}}_i).
\end{equation}

Operation Accuracy is calculated in the same manner for the supported
operations: lookup, recommend, compare, plan, and explain. Intent and operation
are evaluated separately because they represent the domain and requested
action, respectively.

\subsection{Constraint and Facet F1}

Explicit constraints and retrieval facets are normalized into key--value
items. For a predicted item set $P$ and reference set $R$:

\begin{equation}
\operatorname{Precision}=\frac{|P\cap R|}{|P|},\qquad
\operatorname{Recall}=\frac{|P\cap R|}{|R|},
\end{equation}

\begin{equation}
F_1=\frac{2\cdot\operatorname{Precision}\cdot\operatorname{Recall}}
{\operatorname{Precision}+\operatorname{Recall}}.
\end{equation}

Query Constraint F1 uses explicit hard requirements, while Retrieval Facet F1
uses normalized semantic signals. The empty-set cases are handled
deterministically so that correctly predicting no items receives full credit.

\section{Retrieval Metrics}

The independent judge grades unique candidate chunks from 0 to 3:

\begin{itemize}
    \item 0: unrelated to the request;
    \item 1: provides general background only;
    \item 2: useful for part of the request; and
    \item 3: directly useful for answering the request.
\end{itemize}

\subsection{Precision@5}

A chunk is considered relevant for Precision@5 when its grade is at least 2:

\begin{equation}
P@5=\frac{1}{5}\sum_{i=1}^{5}\mathbb{I}(rel_i\geq2).
\end{equation}

\subsection{nDCG@5}

The discounted cumulative gain at rank five is:

\begin{equation}
DCG@5=\sum_{i=1}^{5}\frac{2^{rel_i}-1}{\log_2(i+1)}.
\end{equation}

The ideal DCG is obtained by sorting the same grades in descending order. The
normalized score is:

\begin{equation}
nDCG@5=\frac{DCG@5}{IDCG@5}.
\end{equation}

Recall@5 is not reported because the complete set of relevant chunks in the
corpus is not exhaustively annotated. Evidence coverage is retained as a
pipeline diagnostic and is not presented as corpus-level recall.

\section{Final-Answer Metrics}

The judge assigns scores from 1 to 5 for:

\begin{itemize}
    \item \textbf{Correctness}: accuracy of claims and recommendations relative
    to the request and evidence;
    \item \textbf{Faithfulness}: degree to which factual claims are supported
    by selected evidence;
    \item \textbf{Personalization Adherence}: appropriate application of
    relevant preferences, memories, and constraints; and
    \item \textbf{Completeness}: coverage of supported request aspects and clear
    acknowledgment of unsupported aspects.
\end{itemize}

The judge receives the evidence in the exact order supplied to answer
generation and uses temporary evidence identifiers. Retrieval candidates also
use temporary identifiers that are mapped back to corpus UUIDs after judging,
reducing identifier transcription errors.

\section{Additional Diagnostics}

The experiment reports but does not merge the following diagnostics into one
quality score:

\begin{itemize}
    \item coverage ratio before and after recovery;
    \item complete, partial, and insufficient readiness counts;
    \item number of recovery searches and new candidates;
    \item newly covered requirements;
    \item filter relaxations;
    \item likely corpus, reranking, or selected-evidence failure causes;
    \item parser fallback frequency;
    \item stage and end-to-end latency;
    \item token usage; and
    \item estimated cost.
\end{itemize}

These fields help explain why a case obtained its final scores. For example, a
low completeness score may be expected when readiness is partial because the
corpus lacks a required schedule, whereas low faithfulness with complete
readiness suggests a generation problem.

\section{Human Validation and Threats to Validity}

Human review is performed at three points. First, generated user profiles and
queries are checked for coherence, diversity, and compatibility with the
Vietnam travel scope. Second, reference intent, operation, constraints, and
facets are corrected where the synthetic annotation conflicts with the agreed
taxonomy. Third, a representative sample of relevance grades and final-answer
scores is inspected, with priority given to disagreements, extreme values,
fallback cases, and corpus-gap diagnoses.

Internal validity is threatened by stochastic LLM behavior and evaluator
variance. Versioned prompts, schemas, temperatures, model identifiers, and
resumable artifacts reduce uncontrolled variation. Construct validity is
limited because numerical judge scores approximate concepts such as
faithfulness and personalization. Clear rubrics and human inspection reduce but
do not eliminate this limitation.

External validity is limited by synthetic profiles, one national tourism scope,
one corpus snapshot, and 100 queries. The results therefore describe the tested
system configuration rather than every travel domain or real user. The official
100-query set is kept separate from small development smoke runs to reduce
overfitting through repeated manual tuning.

\section{Reproducibility}

Each run records the pipeline version, source-control commit when available,
corpus fingerprint, embedding and reranker models, LLM roles, pricing snapshot,
prompt hashes, timestamps, and random seed. Completed cases are saved
incrementally. The judge can resume completed cases only when the run identity,
pipeline version, judge model, and judge schema remain compatible.

The commands required to reproduce the official run are provided in
Appendix~\ref{app:commands}.

